% !TeX root = ../papers/nc-p3.tex
\section{階数15の同一倍率上の添字保持同型類}
\label{sec:rank-fifteen-root-exchange}
\begin{theorem}\label{thm:rank-fifteen-root-exchange}
標数零の幾何点で、正の固定曲線の通常倍率 $(z,\lambda)$ を固定する。
$zk^2\ne1$ とする。分岐式 $\delta=(1-zk^2)^2-4zm^2$ が非零なら、
二表示の八窓は添字保持同型でちょうど二類に分かれ、各類の最小周期は2である。
$\delta=0$ なら二表示の四窓は一類になる。ここで数えるのは同じ倍率の上の窓である。
\end{theorem}
\begin{proof}
$r^2=z$, $s^2=r$, $d^2=\lambda$, $i^2=-1$ と取る。
$n$ を定理\ref{thm:rank-fifteen-covers}の四次式の一根とすると、
$N=rm/n$ により全根は $n,-n,N,-N$ である。
$\lambda'=1/(z\lambda)$ としても $k,m$ は変わらない。
$H,H'$ を第一表示の $(z,\lambda,n),(z,\lambda',N)$ の規格化テンソルとする。
前節の出力・入力の行列 $A_0,A_1$ を用いて
\[
 G=D(1,1,(ds)^{-1},(ds)^{-1}),\quad
 G_0=G(I-iA_0/r),\qquad G_1=G(I+iA_1)
\]
と置く。$A_0^2=zI$, $A_1^2=I$ なので両行列は可逆である。すると
\[
 H'(G_0\otimes G_1)=(G_0\otimes G_1)H .
\]
\path{computations/verify_rank_fifteen_root_exchange.py} は元のGramブロックで各成分を展開する。
四変数 $s,d,n,i$ の有理式で224成分の差が恒等的に零で、残る32成分の分子は
$i^2+1$ で割り切れる。分母には $i$ を使わず、二つのGram逆行列は
定理\ref{thm:rank-fifteen-covers}により両枠で通常点に極を持たない。
この共役は始点を保つ。定理\ref{thm:rank-fifteen-endpoint-orbits}と合成すると
同じ倍率で第一表示の $n$ と第二表示の $N$ が同型になる。
また $n\mapsto-n$ は二位置移動であり、定理\ref{thm:rank-fifteen-word-symmetries}で
同じ窓はこの移動と添字保持同型である。従って高々二類で、$N=\pm n$ となる分岐では一類である。

二類を区別するため、$B(S)$ を命題\ref{prop:rank-fifteen-oriented-cubic}の $C$ の
分子とし、$S_1=T^*T$, $S_0=TT^*$ と置く。階数15ではトレース零空間上の $T$ は可逆なので
\[
 J=\frac{(B(S_1)-B(S_0))^5}{\det(S_1|_{\mathfrak{sl}_4})}
\]
は全通常窓で定義される。$T$ のスカラー倍に関する分子・分母の次数はともに30で、
両方とも独立な基底共役で不変である。語の反転と一位置移動はいずれも二つの
行列空間を交換する。実際、各位置をその反対位置と核形式で対にし、
$H$ の行・列の組を $(0,2),(1,3)$ とすれば、反転後の組は $(3,1),(2,0)$、
一位置移動後は $(1,3),(2,0)$ となる。対の順序を逆にすると、
核形式での基底同一視を除いて行列を転置する。転置もトレース形式と
$d(X,Y,Z)=\operatorname{tr}((XY+YX)Z)$ を保つので、いずれの操作でも二つの $B$ は交換する。
$S_0$ と $S_1$ は共役なので分母は等しく、$J$ の符号は反転する。

最後に $B(S_1)-B(S_0)$ の零点を調べる。上記コードで $H$ の四文字別の $T$ を
共通分母倍して多項式行列にし、トレース随伴を取り、二つの三次縮約を引く。
$P(n)$ の分子で割る前に最高係数の単元の冪を掛ける。剰余は $n$ に関して二次であり、
$n,z,\lambda$ の共通単項式を除いて $P$ の分子との終結式を取ると、非零定数を除いて
\[
 \begin{aligned}
 &\operatorname{num}(\delta)^2 (z-1)^4 z^{28}
 (z\lambda^2-1)^{52}\\
 &\hspace{20mm}{}\cdot\lambda^{56}(\lambda-1)^{60}(z\lambda-1)^{60}
 \end{aligned}
\]
となる。分岐因子以外は通常倍率で単元であり、$n$ も単元である。
従って非分岐点では $J\ne0$ で、同じ $n$ の反対表示も一位置移動も元の窓と区別される。
二位置移動の同型と合わせて最小周期も従う。
\end{proof}

\begin{corollary}
定理\ref{thm:rank-fifteen-coordinate-collisions}の楕円成分の $U=3$ では、
同じ $s$ の二点が同型になり、二つの $s$ は相異なる類を与える。
各有理成分の $U=-1$ ではその成分の二点が一類になる。
\end{corollary}
\begin{proof}
楕円成分の $n=1$ では $P(1)=0$ から $N^2=U-2$、有理成分の $m=ik,n=i$ では
$N^2=U$ である。指定水準では $N=\pm n$ なので、第一表示内で
$\lambda\mapsto1/(z\lambda)$ が同型を与える。楕円上では $t=z$ を固定して
$Y$ の符号を交換するので $s=t+1/t$ を保ち $v$ の符号を交換する。
$s^2-28s+4=0$ の異なる二根の間の同型は前掲定理の $s\ne u$ の消去で除かれる。
有理成分では $t=\lambda/i$ が $(t+1)/(t-1)$ に移り、$s=t-1/t$ は $4/s$ に移る。
これは $U=-1$ の水準式の異なる二根を交換する。共役成分でも同じ計算が成り立つ。
\end{proof}
